%==============================================================================
%  Bhargavi Gudisa — Résumé
%
%  Compile with pdflatex (Overleaf: "pdfLaTeX"). One column, selectable text,
%  no graphics and no text inside tables that carries meaning, so applicant
%  tracking systems parse it the same way a human reads it.
%
%  Everything you are likely to edit is in the PERSONAL DETAILS block below.
%  Lines marked  % TODO  need a real value before you send this out.
%
%  FITTING: this should land on two pages. If it spills onto a third, cut in
%  this order — the goal is two full pages, never two and a bit:
%    1. the Permit Request Tracker project (weakest of the four, no link yet)
%    2. the last two Wipro bullets (oldest role, least relevant to cloud work)
%    3. the "Practice" skills row
%  If instead you want one page for a competitive SWE req, keep Lumin's first
%  four bullets, FAU's first three, one line for Wipro, and two projects.
%==============================================================================
\documentclass[letterpaper,10pt]{article}

\usepackage[T1]{fontenc}
\usepackage[utf8]{inputenc}
\usepackage{charter}                 % readable serif; ships with every TeX distro
\usepackage{microtype}
\usepackage[top=0.5in,bottom=0.5in,left=0.55in,right=0.55in]{geometry}
\usepackage{titlesec}
\usepackage{enumitem}
\usepackage{xcolor}
\usepackage{tabularx}
\usepackage[colorlinks=true,urlcolor=linkblue,linkcolor=linkblue]{hyperref}

\definecolor{linkblue}{HTML}{1A4F8A}
\definecolor{rulegray}{HTML}{333333}

\pagestyle{empty}
\setlength{\parindent}{0pt}
\raggedbottom
\raggedright

%--- Section headings ---------------------------------------------------------
\titleformat{\section}
  {\normalsize\bfseries\scshape\color{rulegray}}
  {}{0em}{}
  [\vspace{-7pt}\color{rulegray}\rule{\linewidth}{0.7pt}\vspace{-3pt}]
\titlespacing*{\section}{0pt}{9pt}{5pt}

%--- Bullets ------------------------------------------------------------------
\newenvironment{bullets}
  {\begin{itemize}[leftmargin=1.15em,labelsep=0.45em,itemsep=1.3pt,
                   topsep=2.5pt,parsep=0pt,partopsep=0pt,label=\textbullet]}
  {\end{itemize}\vspace{1pt}}

%--- One experience / education entry: {left}{right}{sub-left}{sub-right} ------
\newcommand{\entry}[4]{%
  \begin{tabularx}{\linewidth}{@{}X r@{}}
    \textbf{#1} & #2 \\
    \textit{#3} & \textit{#4} \\
  \end{tabularx}%
  \vspace{-3pt}
}

%--- A project line: {name}{stack}{links} -------------------------------------
\newcommand{\project}[3]{%
  \begin{tabularx}{\linewidth}{@{}X r@{}}
    \textbf{#1}~~\textbar{}~~\textit{\small #2} & \small #3 \\
  \end{tabularx}%
  \vspace{-3pt}
}

%==============================================================================
%  PERSONAL DETAILS — edit here
%==============================================================================
% NOTE: each link is defined as a complete \href. Do not split the URL into its
% own \newcommand — hyperref reads the URL argument with special catcodes and
% does not reliably expand macros inside it.
\newcommand{\myname}{Bhargavi Gudisa}
% No dash between the two halves. For a wider gap instead of a single space,
% put \quad between them.
\newcommand{\mytitle}{Software Engineer  Cloud Infrastructure \& AI Quality}

\newcommand{\myphone}{\href{tel:+15619870286}{(561) 987-0286}}
\newcommand{\myemail}{\href{mailto:bhargavi4044g@gmail.com}{bhargavi4044g@gmail.com}}

\newcommand{\mylocation}{USA}

% Short labels rather than printed URLs. One tradeoff to know about: an ATS that
% discards hyperlink targets keeps only the word "GitHub", where a printed URL
% would have survived as text. Most modern parsers read the target, and nothing
% is lost on a human reader. To go back to full URLs, swap in the commented
% versions below.
\newcommand{\mysite}{\href{https://bhargavigudisa.github.io}{Portfolio}}
\newcommand{\mylinkedin}{\href{https://www.linkedin.com/in/bhargavi-gudisa-a2b05021b/}{LinkedIn}}
\newcommand{\mygithub}{\href{https://github.com/BhargaviGudisa}{GitHub}}

% \newcommand{\mysite}{\href{https://bhargavigudisa.github.io}{bhargavigudisa.github.io}}
% \newcommand{\mylinkedin}{\href{https://www.linkedin.com/in/bhargavi-gudisa-a2b05021b/}{linkedin.com/in/bhargavi-gudisa-a2b05021b}}
% \newcommand{\mygithub}{\href{https://github.com/BhargaviGudisa}{github.com/BhargaviGudisa}}

\newcommand{\sep}{~~\textbullet~~}
\newcommand{\pipe}{~~\textcolor{rulegray}{\textbar}~~}

%==============================================================================
\begin{document}

%--- HEADER -------------------------------------------------------------------
\begin{center}
  {\LARGE\bfseries\myname}\\[3pt]
  {\small\mytitle}\\[5pt]
  {\small
    \myphone\sep\myemail\sep\mylocation\\[3pt]
    \mygithub\pipe\mylinkedin\pipe\mysite
  }
\end{center}
\vspace{-4pt}

%--- SUMMARY ------------------------------------------------------------------
\section{Summary}

Software engineer with 5 years across the US and India building cloud
infrastructure and the quality systems that run on top of it. Provisions GCP and
Azure environments with Terraform, packages Python and Bash pipelines into Docker
on Linux, and runs AI response evaluation with human-in-the-loop review and
root-cause routing so each failure reaches the one team that can fix it. M.S. in
Computer and Information Sciences, Florida Atlantic University (GPA 3.86).

%--- SKILLS -------------------------------------------------------------------
\section{Skills}

\begin{tabularx}{\linewidth}{@{}>{\bfseries}l X@{}}
  Infrastructure \& Cloud~ & Terraform, Docker, Infrastructure as Code (IaC),
    GCP (Cloud Functions, Cloud Storage, Compute Engine), Microsoft Azure
    (Data Factory, Synapse, App Service), Linux, Git, GitHub \\[2.5pt]
  Languages & Python, Bash, SQL (T-SQL), C\#, JavaScript \\[2.5pt]
  Data \& Reporting & Microsoft SQL Server, Azure SQL Database, Oracle PL/SQL,
    ETL, pandas, relational data modeling, indexing and query tuning,
    Power BI, Tableau, DAX, SSRS \\[2.5pt]
  AI \& Integration & AI quality evaluation, human-in-the-loop review, Microsoft
    Copilot Studio, REST APIs, Workato, Make, n8n, Power Automate,
    SharePoint Online \\[2.5pt]
  Practice & SDLC, Agile, design and code reviews, test planning, data migration
    and conversion, technical documentation, Jira, Asana \\
\end{tabularx}

%--- EXPERIENCE ---------------------------------------------------------------
\section{Experience}

\entry{Software Engineer}{May 2025 \textendash{} Present}
      {Lumin Innovations}{Remote, United States}  % TODO confirm location
\begin{bullets}
  \item Built a Python and Bash ingestion application, packaged it in Docker, and
        provisioned its GCP storage and compute with Terraform, so every Linux
        environment runs the same version-controlled build.
  \item Set up knowledge-base governance with an owner per content area and
        Jira-tracked reviews each release; automated Bash audits route 200+ stale
        entries a month to their owners and cut staleness by 45\%.
  \item Built Python and Bash tooling that analyzed 5,000+ AI-generated chat
        responses, flagging quality gaps at 98\% accuracy in remediation reports
        engineers used to prioritize fixes.
  \item Designed a root-cause process separating content gaps from prompt-handling
        failures, so engineers check one cause per failure instead of both;
        debugging time fell 30\%.
  \item Evaluated, configured, and integrated SaaS automation platforms (Workato,
        Make, n8n) into content and training workflows, replacing manual handoffs
        between teams with automated triggers.
  \item Reported AI accuracy and adoption to executives monthly through Tableau
        dashboards and set cross-functional priorities in Asana; response
        precision rose 25\%.
\end{bullets}

\entry{Systems Analyst}{Aug 2024 \textendash{} May 2025}
      {Florida Atlantic University}{Boca Raton, FL}
\begin{bullets}
  \item Designed a layered GCP ingestion system with Cloud Functions and Terraform
        that separated collection, processing, and storage, making new data
        sources quick to onboard.
  \item Deployed the pipeline in Docker on Linux virtual machines and optimized its
        Bash metadata extraction scripts, keeping development and production
        consistent and cutting environment-specific failures.
  \item Defined quality criteria for AI system responses and built a 500+ case
        edge-case test dataset, giving each release a repeatable, traceable test
        baseline.
  \item Worked with documentation owners in Jira to fix low-confidence knowledge
        articles, raising search success 18\%; content audits cut outdated entries
        40\% and saved about 15 hours of manual review a week.
  \item Ran cross-team workshops to align system output with institutional policy,
        reducing policy-related escalations by 22\%.
\end{bullets}

\entry{Project Engineer}{Jan 2021 \textendash{} Oct 2023}
      {Wipro Limited}{India}  % TODO add the city (Hyderabad? Bangalore?)
\begin{bullets}
  \item Containerized the QA framework with Docker and extended Terraform
        configurations for GCP, cutting manual setup as test volume grew and
        speeding developer onboarding.
  \item Migrated data processing scripts to Bash on Linux servers, making them
        portable across environments and laying the base for later GCP pipelines.
  \item Ran automated test execution and pre-release validation across the SDLC,
        and joined engineers in design and code reviews to check changes against
        quality specifications.
  \item Reviewed chatbot replies as the human check on AI output, flagging
        inaccurate information to protect data integrity across the conversational
        platform.
  \item Coordinated sprint planning in Jira and Asana across global teams in
        multiple time zones, and built Tableau dashboards that gave leadership one
        view of program performance.
\end{bullets}

%--- PROJECTS -----------------------------------------------------------------
\section{Projects}

% The live demo goes first on purpose: it is the one a reader can verify in
% ten seconds without taking your word for anything.
\project{Review Sentiment \& Competitor Radar}
        {Python, Streamlit, Plotly, pandas, SQLite, VADER}
        {\href{https://github.com/BhargaviGudisa/review-sentiment-radar}{Code}~~\textbar{}~~\href{https://review-sentiment-radar-iyxpjovmbfqcbyeuzk3ifs.streamlit.app/}{Live demo}}
\begin{bullets}
  \item Built an aspect-based sentiment pipeline over the Yelp Open Dataset
        (6.99M reviews) that scores a business against five nearby competitors
        across six service aspects, rather than collapsing reviews to one number.
  \item Scored sentiment per sentence so a single review can praise the coffee and
        damn the wait without the two cancelling out; chose VADER and keyword
        matching over an LLM for cost, speed, and explainability.
  \item Shipped a Streamlit dashboard with radar comparison and automated gap
        analysis, surfacing the three aspects where the business trails most.
\end{bullets}

\project{Record Retention Advisor}
        {Microsoft Copilot Studio, Microsoft 365}
        {Town of Morrisville, NC}
\begin{bullets}
  \item Built a Copilot Studio agent that answers Town staff questions on how long
        to keep municipal records and how to handle them when retention ends.
  \item Grounded the agent in the Town's own retention schedule through knowledge
        sources and topics, replacing manual schedule lookups with a conversational
        interface inside Microsoft 365.
  \item Made every answer cite its schedule entry, so staff verify against the
        source before acting \textemdash{} destroying a record early is not a
        recoverable mistake.
\end{bullets}

% TODO: publish this repo, then replace the {} below with:
%   {\href{https://github.com/BhargaviGudisa/coffee-quality-pipeline}{Code}}
% Until it exists, the "published to GitHub" claim in the last bullet is not
% checkable, which is worse than not making it.
\project{Coffee Sensory Quality Data Pipeline}
        {Python, pandas, Azure SQL Database, Power BI}
        {}
\begin{bullets}
  \item Built a Python pipeline that cleans the public Coffee Quality Institute
        dataset (1,300+ Arabica cupping reviews), applies named quality rules, and
        loads it into Azure SQL Database.
  \item Logged every rejected record with the rule it failed, so each loaded score
        traces back to its source row and no record drops silently.
  \item Published the data lineage and a data dictionary so other analysts can
        reuse the model rather than re-derive it.
\end{bullets}

% TODO: same as above —
%   {\href{https://github.com/BhargaviGudisa/permit-request-tracker}{Code}}
\project{Permit Request Tracker}
        {C\#, ASP.NET Core MVC, Entity Framework Core, SQL Server, Azure App Service}
        {}
\begin{bullets}
  \item Built a web application where staff log, assign, and close permit requests,
        modelled with Entity Framework Core over a normalized SQL Server schema.
  \item Wrote a T-SQL conversion script migrating legacy CSV records, coercing dates
        safely and refusing unparseable rows instead of inserting nulls.
  \item Chose indexes from before-and-after execution plans rather than intuition,
        and deployed to Azure App Service.
\end{bullets}

%--- EDUCATION ----------------------------------------------------------------
\section{Education}

% TODO: add the start date. Right now the résumé shows Wipro ending Oct 2023 and
% FAU starting Aug 2024, with nothing in between — a reader sees a 10-month gap.
% The M.S. start date closes it. Replace "Jan 2024" with your actual start.
\entry{M.S., Computer and Information Sciences \textemdash{} GPA 3.86}
      {Jan 2024 \textendash{} Dec 2025}
      {Florida Atlantic University}{Boca Raton, FL}
\vspace{2pt}

Relevant coursework: Cloud Computing, Database Systems, Big Data Analytics,
Machine Learning, Data Mining, Data Visualization.
\vspace{4pt}

%--- CERTIFICATIONS -----------------------------------------------------------
\section{Certifications \& Publications}

\begin{bullets}
  \item \textbf{Microsoft Certified: Azure Fundamentals (AZ-900)} \textemdash{}
        Microsoft.
        % TODO: paste your Credly / Microsoft Learn verification URL here, e.g.
        % \href{https://learn.microsoft.com/api/credentials/share/...}{Verify}
  \item \textbf{Programming Essentials in Python} \textemdash{} Statement of
        Achievement, Cisco Networking Academy.  % TODO confirm the issuer
  \item \textbf{Hybrid model of a passive aggressive classifier machine learning
        approach to fake news detection} \textemdash{} Springer.
        % TODO: add the publication year and a DOI link. A cited paper a reader
        % cannot find does less for you than one they can open:
        % \href{https://doi.org/XX.XXXX/XXXXXX}{DOI}
\end{bullets}

\end{document}
